\documentclass{article}

\usepackage[final]{neurips_2020}

\usepackage[utf8]{inputenc}
\usepackage[T1]{fontenc}
\usepackage{hyperref}
\usepackage{url}
\usepackage{booktabs}
\usepackage{amsfonts}
\usepackage{nicefrac}
\usepackage{microtype}
\usepackage{graphicx}
\usepackage{xcolor}
\usepackage{lipsum}

\newcommand{\note}[1]{\textcolor{blue}{{#1}}}

\title{
  Title of your project \\
  \vspace{1em}
  \small{\normalfont Korea University COSE461 Final Project}  % Select one and delete the other
}

\author{
  Name \\
   Department of Computer Science \\
   Team 0 \\
   2099012345
  % Examples of more authors
%   \And
%   Name \\
%   Department of Computer Science \\
%   Team 0 \\
%   2099012345
%   \And
%   Name \\
%   Department of Computer Science \\
%   Team 0 \\
%   2099012345
}

\begin{document}

\maketitle

\begin{abstract}
An abstract should concisely (less than 300 words) motivate the problem, describe your aims, describe your contribution, and highlight your main finding(s). 
\end{abstract}

This project report is limited to 8 content pages (excluding reference); i.e. you can submit reports with 1-8 pages. Please make your report brief and concise, so that there is no redundant contents. There is no bonus point on submitting full 8 pages report. 

\section{Introduction}
The introduction explains the problem, why it's difficult, interesting, or important, how and why current methods succeed/fail at the problem, and explains the key ideas of your approach and results. Though an introduction covers similar material as an abstract, the introduction gives more space for motivation, detail, references to existing work, and to capture the reader's interest.

\section{Related Work}
This section helps the reader understand the research context of your work, by providing an overview of existing work in the area.
\begin{itemize}
    \item You might discuss: papers that inspired your approach, papers that you use as baselines, papers proposing alternative approaches to the problem, papers applying your methods to different tasks, etc.
    \item This section shouldn't go into deep detail in any one paper (for example, there probably shouldn't be any equations) – instead it should explain how the papers relate to each other, and how they relate to your work. If you need to go deep into detail in previous researches with equations, discuss it in the next section.
\end{itemize}

\section{Approach}
This section details your approach(es) to the problem. 
For example, this is where you describe the architecture of your neural network(s), and any other key methods or algorithms.
\begin{itemize}
    \item You should be specific when describing your main approaches – you probably want to include equations and figures.
    \item You should also describe your baseline(s). Depending on space constraints, and how standard your baseline is, you might do this in detail, or simply refer the reader to some other paper for the details.
    \item If any part of your approach is original, make it clear (so I can give you credit!). For models and techniques that aren't yours, provide references.
    \item If you’re using any code that you didn't write yourself, make it clear and provide a reference or link. When describing something you coded yourself, make it clear (so I can give you credit!).
    \item As you’re setting up equations, notation, and the like, be sure to agree on a fixed technical vocabulary (that you've defined, or is well-defined in the literature) before writing and use it consistently throughout the report!

\end{itemize}

\section{Experiments}
This section contains the following.

\subsection{Data}
Describe the dataset(s) you are using (provide references). If it's not already clear, make sure the associated task is clearly described.
Being precise about the exact form of the input and output can be very useful for readers attempting to understand your work, especially if you've defined your own task.

\subsection{Evaluation method}
Describe the evaluation metric(s) you use, plus any other details necessary to understand your evaluation.
Some projects will have clear metrics from prior work on given datasets, but we realize that other projects will define their own metrics.
If you're defining your own metrics, be clear as to what you're hoping to measure with each evaluation method (whether quantitative or qualitative, automatic or human-defined!), and how it's defined.

\subsection{Experimental details}
Report how you ran your experiments (e.g. model configurations, learning rate, training time, etc.)

\subsection{Results}
Report the quantitative results that you have found so far. Use a table or plot to compare results and compare against baselines.
\begin{itemize}
    \item Comment on your quantitative results. Are they what you expected? Better than you expected? Worse than you expected? Why do you think that is? What does that tell you about your approach?
\end{itemize}

\section{Analysis}
Your report should include \textit{qualitative evaluation}. That is, try to understand your system (e.g. how it works, when it succeeds and when it fails) by inspecting key characteristics or outputs of your model.

\section{Conclusion}
Summarize the main findings of your project, and what you have learnt. Highlight your achievements, and note the primary limitations of your work. If you like, you can describe avenues for future work.

After you are done writing this report, save it as a .pdf file, and submit through BlackBoard. You don't have to submit a copy per person; submitting once per team is fine.

\bibliographystyle{unsrt}
\bibliography{references}
% you can use .bib file to take care of your references.

\newpage
\appendix

\section{Appendix: Team contributions}
The stuff here does not count towards the 8 page limit. If you are a multi-person team, you can write a brief summary of what each team member did for the project (about 1 or 2 sentences per person). For almost all teams, it will have no effect (i.e. team members all receive same grade), but for teams with considerably unequal contribution, I may investigate and/or give different grades to team members.

\end{document}
